\addsectionframe{3. Lösungsansatz: Architektur und Design}

\begin{itemframe}{Technische Basis}
  \item React 19 und TypeScript bilden die Anwendungslogik.
  \item Vite erzeugt das auslieferbare Frontend; Tailwind CSS und shadcn/ui gestalten die Oberfläche.
  \item Zustand hält wenige globale UI-Zustände, etwa Theme und Seitenleiste.
  \item i18next liefert Texte in Deutsch und Englisch.
\end{itemframe}

\begin{itemframe}{Modulare Architektur}
  \item \texttt{moduleCatalog.tsx} registriert Module, Themen, Routen und Navigation an einer Stelle.
  \item Gemeinsame Seitengerüste liegen in \texttt{components/shared}; fachlicher Inhalt bleibt im jeweiligen Themenordner.
  \item Ganze Seiten werden lazy geladen. Dadurch bleibt die erste Auslieferung kleiner.
  \item Die Trennung verhindert doppelte Registrierung in Router, Header, Sidebar und Modulübersicht.
\end{itemframe}

\begin{itemframe}{Wiederverwendbares Lernmuster}
  \item Jedes Thema verwendet dieselben Bereiche: Überblick, Übungen, Beispiele und Simulation.
  \item \texttt{ModuleTopicPage} steuert Tabs, Aufgabenwahl und den Wechsel zur Simulation.
  \item Themenordner binden nur ihre konkreten Inhalte und Simulatoren ein.
  \item Ergebnis: ein neues Thema folgt einer dokumentierten Ordner- und Registrierungskonvention.
\end{itemframe}

\begin{itemframe}{Designentscheidungen für Simulatoren}
  \item Konfiguration und Ablaufzustand sind getrennt; die Simulation bleibt reproduzierbar.
  \item Schrittsteuerung macht Zustände und Übergänge nachvollziehbar.
  \item Fachmodule nutzen eigene Visualisierungen, etwa Zustandsdiagramme, Wahrheitstabellen oder Bitfolgen.
  \item Import und Export erlauben, einen Arbeitsstand weiterzugeben oder später fortzusetzen.
\end{itemframe}

\begin{itemframe}{Code-Anker für die Befragung}
  \item \texttt{data/moduleCatalog.tsx}: zentrale Registrierung von Modulen, Routen und Navigation.
  \item \texttt{components/shared/ModuleTopicPage/ModuleTopicPage.tsx}: Tabs, Aufgabenwahl und Wechsel zur Simulation.
  \item \texttt{hooks/useSimulationEngine.ts} und \texttt{SimulationWorkspace.tsx}: Ablaufzustand und gemeinsame Simulator-Oberfläche.
  \item Zusätzlich zeigt jedes Teammitglied ein verantwortetes Fachmodul mit einem konkreten Testfall.
\end{itemframe}

\begin{itemframe}{Fachlicher Umfang}
  \item Grundlagen Informatik: Turing-Maschinen, endliche Automaten, reguläre Sprachen, Zahlensysteme, Moore/Mealy, DNF und Kodierungen.
  \item Programmierung 1: Java-Programmierung mit Editor und Ausführung im Browser.
  \item Gemeinsames Ziel aller Module: Eingabe bearbeiten, Ergebnis prüfen und den Lösungsweg sichtbar machen.
\end{itemframe}
